% Source: physical PDF pp. 289--299 (printed pp. 266--276), from the
% Section 29.3 heading through the final paragraph immediately before
% Section 29.4.
% Inventory: Eqs. (29.3.1)--(29.3.38), five unnumbered display groups,
% linked citations [3], [4], and [5], two ordinary footnotes, no figures,
% tables, or centered three-asterisk dividers.
% Corrected the R charge of W, the reality factor in Eq. (29.3.15), the
% instanton condition, and a mistyped cross-reference. Source field-strength
% index slips in Eqs. (29.3.9) and (29.3.32) are normalized here.
% Uncertainties: none.

\subsection{Non-Perturbative Corrections to the Superpotential}

We saw in Section 27.6 that the superpotential in general supersymmetric
theories of gauge and chiral superfields is not renormalized to any finite
order of perturbation theory, so that if supersymmetry is not broken in the
tree approximation then it can only be broken by non-perturbative
corrections to the Wilsonian effective Lagrangian. We now take up a general
analysis of these corrections. These were thoroughly studied in a series of
papers by Affleck, Davis, Dine, and Seiberg in the early
1980s~\hyperref[ch29-ref-3]{[3]}, with special attention to the case of
supersymmetric versions of quantum chromodynamics with arbitrary numbers
of colors and flavors. Here we will present a somewhat simplified analysis
of general supersymmetric gauge theories, based on the more recent
holomorphy arguments of Seiberg~\hyperref[ch29-ref-4]{[4]} already used in
Section 27.7.

To study non-perturbative effects, we shall here again consider a general
renormalizable supersymmetric theory, but now including a possible
\(\theta\)-term in the Lagrangian density
\begin{equation}
  \mathcal{L}
  =
  \left[\Phi^\dagger e^{-V}\Phi\right]_D
  +2\operatorname{Re}\left[f(\Phi)\right]_{\mathcal{F}}
  +\operatorname{Re}
  \left[
    \frac{\tau}{8\pi \ii}
    \sum_AW_A^\alpha W_{A\alpha}
  \right]_{\mathcal{F}}\,,
  \tag{29.3.1}\label{eq:29.3.1}
\end{equation}
where the superpotential \(f(\Phi)\) is a gauge-invariant cubic polynomial
in the left-chiral superfields and \(\tau\) is the parameter
\eqref{eq:27.3.23}:
\begin{equation}
  \tau=\frac{4\pi \ii}{g^2}+\frac{\theta}{2\pi}\,.
  \tag{29.3.2}\label{eq:29.3.2}
\end{equation}

As in Section 27.6, we introduce a pair of gauge-invariant left-chiral
external superfields, now called \(Y\) and \(T\), and replace the
Lagrangian density with
\begin{equation}
  \mathcal{L}^{\#}
  =
  \left[\Phi^\dagger e^{-V}\Phi\right]_D
  +2\operatorname{Re}\left[Yf(\Phi)\right]_{\mathcal{F}}
  +\operatorname{Re}
  \left[
    \frac{T}{8\pi \ii}
    \sum_AW_A^\alpha W_{A\alpha}
  \right]_{\mathcal{F}}\,.
  \tag{29.3.3}\label{eq:29.3.3}
\end{equation}
This becomes the same as \eqref{eq:29.3.1} when we set the spinor and
auxiliary components of \(Y\) and \(T\) equal to zero, and take their
scalar components as \(y=1\) and \(t=\tau\), respectively.
Non-perturbative effects will in general invalidate both of the two
symmetries on which the analysis of Section 27.6 was based. The translation
operation, which in our present notation is \(T\longrightarrow T+\xi\)
with real \(\xi\), is not a symmetry because
\(\sum_A\epsilon_{\mu\nu\rho\sigma}f_A^{\mu\nu}f_A^{\rho\sigma}\) can
have a non-vanishing integral over spacetime. The original \(R\) invariance
(with \(T\) and \(Y\) having \(R\) values 0 and \(+2\)) is not a symmetry
because the anomaly discussed in Chapter 22 gives the \(R\)-current a
non-vanishing divergence. With \(\theta\) and \(\bar\theta\) having
\(R=+1\) and \(R=-1\), respectively, and \(V_A\) and \(\Phi_n\)
\(R\)-neutral, the fermion fields \(\lambda_{A\alpha}\) and
\(\psi_{n\alpha}\)
have \(R=+1\) and \(R=-1\), respectively, so Eq.~\eqref{eq:22.2.26}
here gives
\begin{equation}
  \partial_\mu J_R^\mu
  =
  -\frac{1}{32\pi^2}(C_1-C_2)
  \sum_A\epsilon_{\mu\nu\rho\sigma}
  f_A^{\mu\nu}f_A^{\rho\sigma}\,,
  \tag{29.3.4}\label{eq:29.3.4}
\end{equation}
where \(C_1\) and \(C_2\) are the constants defined in
Eqs.~\eqref{eq:17.5.33} and \eqref{eq:17.5.34}:
\begin{equation}
  \sum_{CD}C_{ACD}C_{BCD}=C_1\delta_{AB}\,,
  \qquad
  \operatorname{Tr}\{t_A t_B\}=C_2\delta_{AB}\,,
  \tag{29.3.5}\label{eq:29.3.5}
\end{equation}
with the trace taken over all species of left-chiral
superfield.\footnote{A factor 32 instead of 16 appears in the denominator
in Eq.~\eqref{eq:29.3.4} because gauginos do not have distinct
antiparticles, and we are now counting antiparticles separately from
particles in taking the trace in Eq.~\eqref{eq:29.3.5}. Also, we are now
adopting the convention, described at the end of Section 27.3, of including
a gauge coupling factor in the gauge fields and not in structure constants
and the matrix generators \(t_A\). The gauge generators are thus normalized
so that for \(t_A\), \(t_B\), and \(t_C\) in the standard \(SU(2)\)
subalgebra of the gauge algebra, the structure constant is
\(C_{ABC}=\epsilon_{ABC}\).} For instance, in the generalized
supersymmetric version of quantum chromodynamics studied in
Reference~\hyperref[ch29-ref-3]{3}, with gauge group \(SU(N_c)\) and
\(N_f\) pairs of left-chiral quark superfields \(Q_n\) and \(\bar Q_n\)
in the defining representation and its complex conjugate, these constants
have values given by Eq.~\eqref{eq:17.5.35} (with \(n_f=2N_f\)) as
\[
  C_1=N_c\,,
  \qquad
  C_2=N_f\,.
\]
Although \(T\) translation and \(R\) invariance are invalidated by
non-perturbative effects, there is a remaining symmetry which is almost as
powerful. Consider a general \(R\) transformation
\begin{equation}
  \theta\longrightarrow e^{\ii\varphi}\theta\,,
  \qquad
  \Phi\longrightarrow\Phi\,,
  \qquad
  V_A\longrightarrow V_A\,,
  \qquad
  Y\longrightarrow e^{2\ii\varphi}Y\,,
  \tag{29.3.6}\label{eq:29.3.6}
\end{equation}
with arbitrary real \(\varphi\). This leaves the \(T\)-independent terms
of the Lagrangian density \eqref{eq:29.3.3} invariant, but according to
Eq.~\eqref{eq:29.3.4}, quantum effects violate this symmetry, just as if
there were a term \(\Delta\mathcal{L}\) in the Lagrangian density with a
transformation
\[
  \Delta\mathcal{L}
  \longrightarrow
  \Delta\mathcal{L}
  -\frac{1}{32\pi^2}(C_1-C_2)
  \sum_A\epsilon_{\mu\nu\rho\sigma}
  f_A^{\mu\nu}f_A^{\rho\sigma}\varphi\,.
\]
Recalling Eq.~\eqref{eq:27.3.18}, this is cancelled if we give \(T\) a
transformation
\begin{equation}
  T\longrightarrow T+(C_1-C_2)\varphi/\pi\,.
  \tag{29.3.7}\label{eq:29.3.7}
\end{equation}
Because \(W_{A\alpha}\) has \(R=1\), the whole theory including
non-perturbative effects is invariant under the combined transformations
\eqref{eq:29.3.6} and \eqref{eq:29.3.7}. In particular, the superfield
\(\exp(2\ii\pi T)\), which for \(T=\tau\) is periodic in \(\theta\), has
\(R=2(C_1-C_2)\).

We again introduce an ultraviolet cut-off, and consider the effective
`Wilsonian' Lagrangian
\begin{equation}
  \begin{aligned}
  \mathcal{L}^{\#}_\lambda
  ={}&
  \left[
    \mathcal{A}_\lambda
    \bigl(\Phi,\Phi^\dagger,V,T,T^\dagger,Y,Y^\dagger,
    \mathcal{D}\cdots\bigr)
  \right]_D
  \\
  &+2\operatorname{Re}
  \left[
    \frac{T}{8\pi \ii}
    \sum_AW_A^\alpha W_{A\alpha}
    +\mathcal{B}_\lambda(\Phi,W,T,Y)
  \right]_{\mathcal{F}}\,,
  \end{aligned}
  \tag{29.3.8}\label{eq:29.3.8}
\end{equation}
with \(\mathcal{A}_\lambda\) and \(\mathcal{B}_\lambda\) both
gauge-invariant functions of the displayed arguments. The term
proportional to \(T\) has been separated from the function
\(\mathcal{B}_\lambda\) in order that the translation
\eqref{eq:29.3.7} of \(T\) should continue to cancel the anomaly in the
\(R\) transformation \eqref{eq:29.3.6}. Invariance under the combined
transformation \eqref{eq:29.3.6}, \eqref{eq:29.3.7} then tells us that
terms in the function \(\mathcal{B}_\lambda\) must be proportional to
powers of \(\exp(2\ii\pi T)\), which have definite \(R\) values.

Furthermore, it is only \emph{positive} powers of \(\exp(2\ii\pi T)\) that
may appear in \(\mathcal{B}_\lambda\). According to
Eq.~\eqref{eq:27.3.24}, it is only instantons with positive winding number
\(\nu\geq0\) that can make contributions to the effective Lagrangian that
are holomorphic in \(T\) rather than \(T^*\), and these give rise to
factors \(\exp(2\ii\pi\nu T)\). More generally, for \(T=\tau\) any power
\(\exp(2ia\pi T)\) will depend on the gauge coupling through a factor
\(\exp(-8\pi^2a/g^2)\), so that \(a\) must be positive in order for
non-perturbative effects to be suppressed for small \(g\). In consequence,
non-perturbative effects enter in \(\mathcal{L}^{\#}_\lambda\) through
operators \(\exp(2\ii\pi aT)\) that have positive-definite, zero, or
negative-definite values of \(R\), depending on whether \(C_1>C_2\),
\(C_1=C_2\), or \(C_1<C_2\). (In the generalized supersymmetric version
of quantum chromodynamics described above, this corresponds to \(N_c>N_f\),
\(N_c=N_f\), and \(N_c<N_f\), respectively.) We shall now consider each
of these cases.

\begin{center}
  \(C_1>C_2\)
\end{center}

Here the powers \(\exp(2ia\pi T)\) with \(a>0\) have positive-definite
values of \(R\), given by Eq.~\eqref{eq:29.3.7} as
\(R=2(C_1-C_2)a\). Lorentz invariance tells us that if any term in
\(\mathcal{B}_\lambda\) contains a factor \(W_{A\alpha}\), then it must
contain at least two of them, so the only ways to construct terms in
\(\mathcal{B}_\lambda\) with \(R=2\) is to have two \(W\)s and no
dependence on \(Y\) or \(T\), or one \(Y\) and no dependence on \(W\) or
\(T\), or one factor of \(\exp(2\ii\pi T/(C_1-C_2))\) and no dependence
on \(W\) or \(Y\):
\begin{equation}
  \begin{aligned}
  \mathcal{B}_\lambda
  ={}&
  Yf_\lambda(\Phi)
  +\sum_{AB}W_A^\alpha W_{B\alpha}
  \ell_{\lambda AB}(\Phi)
  \\
  &+\exp\left(\frac{2\ii\pi T}{C_1-C_2}\right)v_\lambda(\Phi)\,.
  \end{aligned}
  \tag{29.3.9}\label{eq:29.3.9}
\end{equation}
Because \(f_\lambda(\Phi)\) does not depend on \(Y\) or \(T\), it can only
be the tree-approximation superpotential
\begin{equation}
  f_\lambda(\Phi)=f(\Phi)\,,
  \tag{29.3.10}\label{eq:29.3.10}
\end{equation}
just as in perturbation theory. Likewise, because
\(\ell_{\lambda AB}(\Phi)\) does not depend on \(Y\) or \(T\) it must
have equal numbers of \(\Phi\)s and \(\Phi^\dagger\)s, so since it does
not depend on \(\Phi^\dagger\) it cannot depend on \(\Phi\) either. Gauge
invariance then requires (for a simple gauge group) that
\(\ell_{\lambda AB}(\Phi)\) is proportional to \(\delta_{AB}\), and since
it does not depend on \(T\) or \(Y\), the power-counting argument of
Section 27.6 shows that the coefficient of \(\delta_{AB}\) can only be the
one-loop contribution to the running inverse-square Wilsonian gauge
coupling.

\begingroup
\emergencystretch=1em
To be more explicit about the running gauge coupling, recall that in
non-super\-sym\-met\-ric gauge theories with fermions,
Eq.~\eqref{eq:18.7.2}
gives the one-loop renormalization group equation as
\par
\endgroup
\begin{equation}
  \lambda\frac{\dd g_\lambda}{\dd\lambda}=bg_\lambda^3\,,
  \tag{29.3.11}\label{eq:29.3.11}
\end{equation}
with
\begin{equation}
  b=-\frac{1}{4\pi^2}
  \left(\frac{11}{12}C_1-\frac16 C_2\right),
  \tag{29.3.12}\label{eq:29.3.12}
\end{equation}
with the coefficient of \(C_2\) taken as \(-1/6\) rather than \(-1/3\)
because we are now counting the left-chiral states of antifermions
separately from the left-chiral states of particles. As we have seen in
Section 28.2, the effect of gauginos is to multiply the \(C_1\)-term by a
factor \(9/11\), while the effect of the scalar components of the
left-chiral superfields (such as squarks and sleptons) is to multiply the
\(C_2\)-term by a factor \(3/2\), so in supersymmetric theories
Eq.~\eqref{eq:29.3.12} becomes instead
\begin{equation}
  b=-\frac{1}{16\pi^2}(3C_1-C_2)\,.
  \tag{29.3.13}\label{eq:29.3.13}
\end{equation}
The solution of Eq.~\eqref{eq:29.3.11} for the running gauge coupling is
then
\begin{equation}
  g_\lambda^{-2}
  =
  g^{-2}
  +\frac{3C_1-C_2}{8\pi^2}\ln\left(\frac{\lambda}{K}\right),
  \tag{29.3.14}\label{eq:29.3.14}
\end{equation}
where \(K\) is an ultraviolet cut-off, introduced to give meaning to the
otherwise ultraviolet-divergent bare gauge coupling \(g\).

To summarize the results so far, setting \(T=\tau\) and \(Y=1\), the
Wilsonian effective Lagrangian for \(C_1>C_2\) takes the form
\begin{equation}
  \begin{aligned}
  \mathcal{L}^{\#}_\lambda
  ={}&
  \left[
    \mathcal{A}_\lambda
    \bigl(\Phi,\Phi^\dagger,V,\tau,\tau^*,\mathcal{D}\cdots\bigr)
  \right]_D
  +2\operatorname{Re}
  \left[
    \frac{\tau_\lambda}{8\pi \ii}
    \sum_AW_A^\alpha W_{A\alpha}
  \right]_{\mathcal{F}}
  \\
  &+2\operatorname{Re}\left[f(\Phi)\right]_{\mathcal{F}}
  +2\operatorname{Re}\left[
    \exp\left(\frac{2\ii\pi\tau_\lambda}{C_1-C_2}\right)
    v_\lambda(\Phi)
  \right]_{\mathcal{F}}\,,
  \end{aligned}
  \tag{29.3.15}\label{eq:29.3.15}
\end{equation}
where
\begin{equation}
  \tau_\lambda=\frac{4\pi \ii}{g_\lambda^2}+\frac{\theta}{2\pi}\,.
  \tag{29.3.16}\label{eq:29.3.16}
\end{equation}
We have been able to replace \(\tau\) with \(\tau_\lambda\) in the
exponential in the last term in Eq.~\eqref{eq:29.3.15}, because the
difference is a constant times \(\ln\lambda\), which yields a power of
\(\lambda\) that can be absorbed into the definition of \(v_\lambda\).

Non-perturbative effects have now been isolated in the last term in
Eq.~\eqref{eq:29.3.15}. This term can be generated by instantons of winding
number \(\nu>0\) if \(C_1-C_2=1/\nu\). (In general \(C_1-C_2\) is a
rational number. For the generalized supersymmetric version of quantum
chromodynamics, \(C_1-C_2=N_c-N_f\) is an integer, so the condition
\(C_1-C_2=1/\nu\) requires that \(N_c=N_f+1\), and then only \(\nu=1\)
instantons contribute. Detailed calculations~\hyperref[ch29-ref-5]{[5]}
in this model show that instantons actually do make such contributions.)
Whether or not it is instantons that generate the non-perturbative
contribution \(v_\lambda(\Phi)\), we can determine its form by considering
the non-anomalous symmetries of the theory. Since this function is
independent of \(Y\), it can be evaluated as if \(Y=0\), so it shares all
the non-anomalous symmetries of the first term in Eq.~\eqref{eq:29.3.1}.
These include the gauge symmetry itself and a global symmetry under
\(\prod_d SU(n(d))\), where \(d\) labels the different irreducible
representations of the gauge group furnished by the left-chiral
superfields, and \(n(d)\) is the number of times representation \(d\)
occurs. (For instance, in generalized supersymmetric quantum
chromodynamics \(d\) takes two values, labelling the \(N_c\) and
\(\bar N_c\) representations of \(SU(N_c)\), and
\(n(N_c)=n(\bar N_c)=N_f\).) Let us label the \(\Phi\)s as
\(\Phi_{ai}^{(d)}\), where \(a\) is a gauge index, and \(i\) is a `flavor'
index, labelling the \(n(d)\) different \(\Phi\)s that transform under the
gauge group according to the representation \(d\). The only way to
construct a function of the \(\Phi\)s that is invariant under the global
symmetry group \(\prod_d SU(n(d))\) is as a product of \(\Phi\)s, with
the \(n(d)\) flavor indices contracted for each \(d\) with the
antisymmetric \(SU(n(d))\) tensor
\(\epsilon_{i_1\ldots i_{n(d)}}\), and with the gauge indices contracted
with constant tensors of the gauge group. (For instance, for generalized
supersymmetric quantum chromodynamics, \(v_\lambda\) must be a function of
the sole invariant
\[
  D\equiv
  \operatorname{Det}_{ij}\sum_a Q_{ai}\bar Q_{aj}\,,
\]
which is non-zero only for \(N_c\geq N_f\).)

In addition to the anomaly-free \(SU(n(d))\) flavor symmetries, there is
also a \(U_d(1)\) symmetry for each of the irreducible representations
\(d\) furnished by the \(\Phi\)s, with all \(\Phi_{ai}^{(d)}\) for a
given \(d\) undergoing the transformation
\begin{equation}
  \Phi_{ai}^{(d)}
  \longrightarrow
  e^{\ii\varphi_d}\Phi_{ai}^{(d)}\,.
  \tag{29.3.17}\label{eq:29.3.17}
\end{equation}
This symmetry is anomalous, with the effects of the anomaly the same as if
the Lagrangian underwent the change
\begin{equation}
  \mathcal{L}
  \longrightarrow
  \mathcal{L}
  -\sum_d\frac{n(d)C_{2\dd}}{32\pi^2}
  \sum_A\epsilon_{\mu\nu\rho\sigma}
  f_A^{\mu\nu}f_A^{\rho\sigma}\varphi_d\,,
  \tag{29.3.18}\label{eq:29.3.18}
\end{equation}
where \(C_{2\dd}\) is the contribution to \(C_2\) of any one left-chiral
scalar superfield belonging to the irreducible representation \(d\) of the
gauge group. The symmetry is restored if we give \(T\) the transformation
property
\begin{equation}
  T\longrightarrow T+n(d)C_{2\dd}\varphi_d/\pi\,.
  \tag{29.3.19}\label{eq:29.3.19}
\end{equation}
Since \(v_\lambda(\Phi)\) is accompanied in Eq.~\eqref{eq:29.3.9} by a
factor \(\exp(2\ii\pi T/(C_1-C_2))\), which undergoes the transformation
\begin{equation}
  \exp\left(\frac{2\ii\pi T}{C_1-C_2}\right)
  \longrightarrow
  \prod_d
  \exp\left(\frac{+2\ii n(d)C_{2\dd}\varphi_d}{C_1-C_2}\right)
  \exp\left(\frac{2\ii\pi T}{C_1-C_2}\right),
  \tag{29.3.20}\label{eq:29.3.20}
\end{equation}
\begingroup
\emergencystretch=1em
we conclude that \emph{for each representation \(d\) of the gauge group
furnished by the left-chiral scalars, \(v_\lambda(\Phi)\) must be a
homogeneous function of the \(\Phi_{ai}^{(d)}\) of negative order
\(-2n(d)C_{2\dd}/(C_1-C_2)\)}. (For instance, in generalized supersymmetric
quantum chromodynamics we have two irreducible representations of
\(SU(N_c)\), the defining and antidefining representations, each with
\(n(d)=N_f\) and \(C_{2\dd}=1/2\), so \(v_\lambda\) is a homogeneous
function of order \(-N_f/(N_c-N_f)\) in the \(Q\) belonging to the
defining representation and of the same order in the \(\bar Q\) belonging
to the antidefining representation. Thus it must be proportional to
\(D^{-1/(N_c-N_f)}\), where \(D\) is the determinant introduced earlier.)
In general \(C_2=\sum_d n(d)C_{2\dd}\), so \(v_\lambda(\Phi)\) is a
homogeneous function of all the \(\Phi\)s, of order
\(-2C_2/(C_1-C_2)\).
\par
\endgroup

This result satisfies an important consistency check. Recall from
Section 27.4 that any superpotential has dimensionality \(+3\) (counting
powers of mass, with \(\hbar=c=1\)), while the scalar superfields \(\Phi\)
like ordinary scalar fields have dimensionality \(+1\), so the
\(\Phi\)-dependent part of \(v_\lambda\) must appear with a coefficient of
dimensionality
\[
  3+\frac{2C_2}{C_1-C_2}
  =
  \frac{3C_1-C_2}{C_1-C_2}\,.
\]
This coefficient does not depend on the gauge coupling or on any of the
couplings or masses in the superpotential, and because we have replaced
the bare coupling \(g\) with \(g_\lambda\) in the last term of
Eq.~\eqref{eq:29.3.15} it cannot depend on the ultraviolet cut-off \(K\)
used to define \(g\) either, so it can only depend on \(\lambda\).
Therefore
\begin{equation}
  v_\lambda(\Phi)
  =
  \lambda^{(3C_1-C_2)/(C_1-C_2)}H(\Phi)\,,
  \tag{29.3.21}\label{eq:29.3.21}
\end{equation}
where \(H(\Phi)\) is a homogeneous function of order
\(-2n(d)C_{2\dd}/(C_1-C_2)\) in the \(\Phi\)s that belong to each
representation \(d\) of the gauge group and is independent of any
parameters of the theory. We can rewrite Eq.~\eqref{eq:29.3.14} in the
form
\begin{equation}
  \tau_\lambda
  =
  \ii\frac{3C_1-C_2}{2\pi}
  \ln\left(\frac{\lambda}{\Lambda}\right)
  +\frac{\theta}{2\pi}\,,
  \tag{29.3.22}\label{eq:29.3.22}
\end{equation}
where \(\Lambda\) is an energy parameter that characterizes the running
gauge coupling, like the \(\Lambda\simeq200\ \mathrm{MeV}\) of quantum
chromodynamics. Thus the last term in Eq.~\eqref{eq:29.3.15} is
\begin{equation}
  \exp\left(\frac{2\ii\pi\tau_\lambda}{C_1-C_2}\right)
  v_\lambda(\Phi)
  =
  \exp\left(\frac{\ii\theta}{C_1-C_2}\right)
  \Lambda^{(3C_1-C_2)/(C_1-C_2)}H(\Phi)\,.
  \tag{29.3.23}\label{eq:29.3.23}
\end{equation}
\emph{The whole effective superpotential, including the non-perturbative
contribution \eqref{eq:29.3.23}, is therefore independent of the floating
cut-off \(\lambda\).}

The function \(H(\Phi)\) is homogeneous and of negative order in \(\Phi\),
so in the absence of a bare superpotential the potential is
positive-definite at finite values of the scalar fields and vanishes only
at infinite field values. In such a theory, there is no stable vacuum
state, and the question of supersymmetry breaking is moot. The vacuum may
be stabilized by adding a suitable bare superpotential. For instance, in
generalized supersymmetric quantum chromodynamics with \(N_f<N_c\) the
only renormalizable bare superpotential is a sum of mass terms
\begin{equation}
  f(Q,\bar Q)=\sum_{ija}m_{ij}\bar Q_{ai}Q_{aj}\,.
  \tag{29.3.24}\label{eq:29.3.24}
\end{equation}
To seek a supersymmetric vacuum state, we need first of all to find what
scalar components \(q_{ai}\) and \(\bar q_{ai}\) satisfy
Eq.~\eqref{eq:27.4.11}, which here reads
\begin{equation}
  \sum_{abi}q_{ai}^*(t_A)_{ab}q_{bi}
  -
  \sum_{abi}\bar q_{ai}^*(t_A)_{ab}\bar q_{bi}
  =0\,,
  \tag{29.3.25}\label{eq:29.3.25}
\end{equation}
for all generators \(t_A\) of \(SU(N_c)\). The gauge interactions (but
not the superpotential) are invariant under simultaneous \(SU(N_c)\)
transformations on the color indices \(a\) of both \(q_{ai}\) and
\(\bar q_{ai}\); under independent \(SU(N_f)\) and
\(\overline{SU(N_f)}\) transformations on the flavor indices \(i\) of
\(q_{ai}\) and \(\bar q_{ai}\), respectively, and under a \(U(1)\)
transformation of both \(q_{ai}\) and \(\bar q_{ai}\) by opposite phases.
Using these symmetries, it is possible to put the general solution of these
conditions in the form
\begin{equation}
  q_{ai}=\bar q_{ai}
  =
  \begin{cases}
    u_i\delta_{ai}, & a\leq N_f\,,\\
    0, & a>N_f\,,
  \end{cases}
  \tag{29.3.26}\label{eq:29.3.26}
\end{equation}
where the \(u_i\) are complex numbers with the same phase. (Here is the
proof. The \(SU(N_c)\) generators \(t_A\) span the space of all traceless
Hermitian matrices, so Eq.~\eqref{eq:29.3.25} is equivalent to the
requirement that
\begin{equation}
  \sum_i q_{ai}^*q_{bi}
  -
  \sum_i\bar q_{ai}^*\bar q_{bi}
  =
  k\delta_{ab}\,,
  \tag{29.3.27}\label{eq:29.3.27}
\end{equation}
for some constant \(k\). By a combined color and flavor transformation
\(q\longrightarrow UqV\) with \(U\) and \(V\) unitary and unimodular, we
can put the matrix \(q\) in the diagonal form \eqref{eq:29.3.26}, and by a
unimodular change of phase of the diagonal elements we can arrange that
they all have the same phase. Then Eq.~\eqref{eq:29.3.27} becomes
\[
  \sum_i\bar q_{ai}^*\bar q_{bi}
  =
  \begin{cases}
    (u_a^2-k)\delta_{ab}, & a\leq N_f\,,\\
    -k\delta_{ab}, & a>N_f\,.
  \end{cases}
\]
The conditions for \(a>N_f\) show that \(k\leq0\). If \(k\) were
non-zero, then the \(\bar q_{ai}\) would furnish \(N_c\) non-zero
orthogonal vectors with \(N_f\) components, which is impossible for
\(N_f<N_c\), so \(k=0\). We can then put the \(\bar q_{ai}\) in the
diagonal form \eqref{eq:29.3.26} by a series of unitary flavor
transformations: first rotate \(\bar q_{1\ii}\) into the 1-direction; then
keeping the 1-direction fixed, rotate in the space perpendicular to this
direction to put \(\bar q_{2\ii}\) in the 2-direction; and so on; and then
perform a unimodular phase transformation so that all diagonal elements
have the same phase. Eq.~\eqref{eq:29.3.27} then shows that the absolute
values of the diagonal elements of \(q_{ai}\) and \(\bar q_{ai}\) are
equal, and by a non-anomalous opposite phase change of \(q_{ai}\) and
\(\bar q_{ai}\) we can arrange that their common phases are equal, as was
to be proved.)

The function \(H\) in Eq.~\eqref{eq:29.3.23} is here
\begin{equation}
  H(q,\bar q)
  =
  \mathcal{J}
  \left[
    \operatorname{Det}_{ij}\sum_a q_{ai}\bar q_{aj}
  \right]^{-1/(N_c-N_f)}
  =
  \mathcal{J}
  \left[\prod_i u_i\right]^{-2/(N_c-N_f)}\,,
  \tag{29.3.28}\label{eq:29.3.28}
\end{equation}
where \(\mathcal{J}\) is a purely numerical constant. (Detailed
calculations show that \(\mathcal{J}=2\) for \(N_f=2\) and \(N_c=3\).)
Adding the terms \eqref{eq:29.3.23} and \eqref{eq:29.3.24}, the complete
effective superpotential is now
\begin{equation}
  f_{\mathrm{total}}(q,\bar q)
  =
  \mathcal{K}
  \left[\prod_i u_i\right]^{-2/(N_c-N_f)}
  +\sum_i m_i u_i^2\,,
  \tag{29.3.29}\label{eq:29.3.29}
\end{equation}
where
\begin{equation}
  \mathcal{K}
  \equiv
  \mathcal{J}
  \exp\left(\frac{\ii\theta}{N_c-N_f}\right)
  \Lambda^{(3N_c-N_f)/(N_c-N_f)}
  \tag{29.3.30}\label{eq:29.3.30}
\end{equation}
and the \(m_i\) are the diagonal elements of the mass matrix that results
when the original mass matrix is subjected to the
\(SU(N_f)\times\overline{SU(N_f)}\) transformation used to put
the scalars in the form \eqref{eq:29.3.26}. The condition
\eqref{eq:27.4.10}, that \(f_{\mathrm{total}}(q,\bar q)\) should be
stationary, has the solution
\begin{equation}
  u_i^2
  =
  \frac{1}{m_i}
  \left(\frac{\mathcal{K}}{N_c-N_f}\right)^{1-N_f/N_c}
  \left(\prod_j m_j\right)^{-1/N_c}.
  \tag{29.3.31}\label{eq:29.3.31}
\end{equation}
Because we have put the scalar fields in a basis in which the \(u_i\) have
a common phase, the \(m_i\) must also have a common phase in this basis.
But the common phase of the \(u_i^2\) is not unique---the \(1/N_c\) powers
in Eq.~\eqref{eq:29.3.31} tell us that the solution is undetermined by a
factor \(\exp(2\ii\pi n/N_c)\), with \(n\) an integer ranging from 0 to
\(N_c-1\). (The two signs of \(u_i\) for a given \(u_i^2\) are physically
equivalent, because the whole theory is invariant under a non-anomalous
symmetry with \(q_{ai}\longrightarrow e^{\ii\pi}q_{ai}\) and
\(\bar q_{ai}\longrightarrow e^{-\ii\pi}\bar q_{ai}\).) The fact that there
are \(N_c\) physically inequivalent solutions will turn up again in our
discussion of the Witten index in the next section.

\begin{center}
  \(C_1=C_2\)
\end{center}

This case is of some interest because, as we saw in
Eq.~\eqref{eq:27.9.3}, the simplest \(N=2\) supersymmetric Yang--Mills
theory, when written in terms of \(N=1\) superfields, contains a single
left-chiral superfield in the adjoint representation, for
which\footnote{Note that here \(C_2\) refers to the representation of the
gauge group furnished by the chiral superfields, so it is the same as the
quantity \(C_2^b\) in Section 27.9, which refers to the representation
furnished by complex scalars, but half of \(C_2^f\), which refers to the
representation furnished by all spinor fields, including the gauginos.}
of course \(C_2=C_1\).

For \(C_2=C_1\) the function \(\exp(2\ii\pi T)\) has \(R=0\), so that its
appearance in \(\mathcal{L}^{\#}_\lambda\) is not restricted by \(R\)
invariance. The general form of the \(\mathcal{F}\)-term in
\(\mathcal{L}^{\#}_\lambda\) is given here by Eq.~\eqref{eq:29.3.9}, but
with the last term absent:
\begin{equation}
  \mathcal{B}_\lambda
  =
  Yf_\lambda\bigl(\Phi,\exp(2\ii\pi T)\bigr)
  +\sum_{AB}W_A^\alpha W_{B\alpha}
  \ell_{\lambda AB}\bigl(\Phi,\exp(2\ii\pi T)\bigr).
  \tag{29.3.32}\label{eq:29.3.32}
\end{equation}
Because \(f_\lambda\) may depend on \(T\), we cannot now conclude that it
is equal to the bare superpotential, but only that it depends linearly on
whatever coupling coefficients and masses appear linearly in the bare
superpotential. In particular, \emph{if there is no superpotential to
begin with, then none is generated by non-perturbative effects.}

To go further, we need to make use of the anomalous chiral symmetry under a
\(U(1)\) transformation of all the \(\Phi_n\). In order for the whole
theory to be invariant under this symmetry, it is necessary to introduce a
separate external left-chiral superfield \(Y_r\) for the terms in the bare
superpotential of order \(r\) in the \(\Phi_n\). Then the theory is
invariant under the combined transformations
\begin{equation}
  \Phi_n\longrightarrow e^{\ii\varphi}\Phi_n\,,
  \qquad
  T\longrightarrow T+C_2\varphi/\pi\,,
  \qquad
  Y_r\longrightarrow e^{-ir\varphi}Y_r\,.
  \tag{29.3.33}\label{eq:29.3.33}
\end{equation}
This symmetry tells us that a term in \(\mathcal{B}_\lambda\) that is of
order \(\mathcal{N}_r\) in the coefficients of the term in the
superpotential of order \(r\) in \(\Phi\) and proportional to
\(\exp(2ai\pi T)\) must be of an order \(\mathcal{N}_\Phi\) in \(\Phi\),
given by
\begin{equation}
  \mathcal{N}_\Phi
  =
  \sum_r r\mathcal{N}_r-2C_2a\,.
  \tag{29.3.34}\label{eq:29.3.34}
\end{equation}

\begingroup
\emergencystretch=1em
The coefficients \(\ell_{\lambda AB}\) of the terms in
\(\mathcal{B}_\lambda\) that are quadratic in \(W\) are shown in
Eq.~\eqref{eq:29.3.32} to be independent of the parameters in the
superpotential for \(C_1=C_2\), so in this case Eq.~\eqref{eq:29.3.34}
becomes
\par
\endgroup
\begin{equation}
  \mathcal{N}_\Phi=-2C_2a\,.
  \tag{29.3.35}\label{eq:29.3.35}
\end{equation}
Thus there can be no terms in \(\ell_{\lambda AB}\) of positive order in
the \(\Phi_r\), and any term in \(\ell_{\lambda AB}\) that is independent
of the \(\Phi_r\) must be independent of \(T\). These
\(\Phi\)-independent terms in \(\ell_{\lambda AB}\) are therefore again
just the one-loop contribution to the running coupling parameter
\(\tau_\lambda\).

The effective superpotential is shown by Eq.~\eqref{eq:29.3.32} to be
linear in the parameters in the superpotential, so all of its terms have
just one \(\mathcal{N}_r=1\) and the others zero. For such a term
Eq.~\eqref{eq:29.3.34} gives
\begin{equation}
  \mathcal{N}_\Phi=r-2C_2a\,.
  \tag{29.3.36}\label{eq:29.3.36}
\end{equation}
A term in the effective superpotential with \(\mathcal{N}_\Phi\) powers of
\(\Phi\) can therefore only arise from terms in the bare superpotential
with \(r\geq\mathcal{N}_\Phi\) powers of \(\Phi\). The terms with
\(r=\mathcal{N}_\Phi\) have \(a=0\), so they are given by the tree
approximation as just the bare superpotential. The only other terms are
non-perturbative corrections with \(r>\mathcal{N}_\Phi\). Such a
non-perturbative term, of a given order in the \(\Phi\)s, can only arise
from terms in the bare superpotential of \emph{higher} order in the
\(\Phi\)s.

\begin{center}
  \(C_1<C_2\)
\end{center}

Here the \(R\) value \(2(C_1-C_2)\) of \(\exp(2\ii\pi T)\) is
\emph{negative}, so positive powers of \(\exp(2\ii\pi T)\) can compensate
for the positive \(R\) values of \(Y\) and \(W_\alpha\), and
\(\mathcal{B}_\lambda\) may therefore contain terms of arbitrary order in
\(Y\) and \(W_\alpha\). Using the chiral symmetry condition
\eqref{eq:29.3.34} and the \(R\) invariance condition
\begin{equation}
  2
  =
  \mathcal{N}_W
  +2\sum_r\mathcal{N}_r
  -2a(C_2-C_1)\,,
  \tag{29.3.37}\label{eq:29.3.37}
\end{equation}
we can, however, set limits on the structure of terms of a given order
\(\mathcal{N}_\Phi\) in the \(\Phi\)s. Eqs.~\eqref{eq:29.3.34} and
\eqref{eq:29.3.37} have a trivial solution with \(\mathcal{N}_r=1\) for
\(r=\mathcal{N}_\Phi\), \(\mathcal{N}_r=0\) for other values of \(r\),
and \(a=\mathcal{N}_W=0\); this solution just represents the presence in
\(\mathcal{B}_\lambda\) of the original bare superpotential, with no
radiative corrections. If there is no superpotential to begin with, then
\(\mathcal{N}_r=0\) for all \(r\), so Eq.~\eqref{eq:29.3.35} does not
allow any terms in the Wilsonian Lagrangian with \(\mathcal{N}_W=0\), and
so \emph{no superpotential can be generated}. (In the supersymmetric
version of quantum chromodynamics this conclusion is usually derived by
noting that there is no possible term in the superpotential that would be
consistent with all symmetries, but as we have now seen, the conclusion is
much more general.) For renormalizable asymptotically free theories there
is a useful limit on the structure of \(\Phi\)-independent terms in
\(\mathcal{B}_\lambda\). The condition of renormalizability tells us that
\(\mathcal{N}_r=0\) for \(r>3\), so by subtracting \(2/3\) of
Eq.~\eqref{eq:29.3.34} from Eq.~\eqref{eq:29.3.37} we find
\begin{equation}
  2
  \geq
  \frac23\mathcal{N}_\Phi
  +\mathcal{N}_W
  +2a\left(C_1-\frac13C_2\right).
  \tag{29.3.38}\label{eq:29.3.38}
\end{equation}
Asymptotic freedom requires that \(3C_1>C_2\), so for
\(\mathcal{N}_\Phi=0\) (or \(\mathcal{N}_\Phi>0\)) each term on the
right-hand side is positive. It follows that we can have no
\(\Phi\)-independent terms of higher than second order in \(W\), and these
terms are also independent of \(T\), so they again represent the one-loop
contribution to the running coupling parameter \(\tau_\lambda\). But for
\(C_1<C_2\) there is no general prohibition against terms of second or
higher order in \(W\) and negative order in \(\Phi\).
